\subsubsection{Evaluation TD3}
\ac{TD3} is evaluated for both rewards, different exploration noises and the presented buffer types, with pink noise and \ac{ER} as default configuration. Firstly, we see that the behavior of the buffers behaves different to the simple environments. As result, the \ac{PER} buffers lead to a suboptimal policy, which results in a tie for the most runs. The graphical evaluation showed that the agent learns to shoot the puck against the boarder, such that it bounces the complete run over the vertical center line. Furthermore, the critic-loss is the lowest of all \ac{TD3} algorithms, such that we assume that the critc overfits to states with initial high \ac{TD}-error. Besides, \ac{BER} has a higher variance than \ac{ER} in the evaluation, indicating that the buffer is more uncertain in its prediction. This leads to a lower overall performance compared to \ac{ER}.

Secondly, white ($\beta = 0$) and blue ($\beta = 0.5$) noise learn slower than the agent with pink noise and have a lower final evaluation reward. Thus, the results of \cite{eberhard2023pink} also apply to the hockey environment by well balancing the exploitation-exploration trade-off.

Finally, \ac{TD3} in particular benefits from the proposed new reward, containing a time penalty, the distance to the goal center and the refusal to start penalty. Thus, \ac{TD3} with new reward converges after$\sim 40000$ environment steps to an almost perfect policy for the proposed opponents, while \ac{TD3} with the original reward converges not in our experiment. Furthermore, the new reward helps \ac{TD3} the beat \ac{SAC} and increases the final average evaluation reward by more than 2.
